\clearpage
\begin{trapbox}
  
    \textbf{\textcolor{CTrap}{易错一：看到 $x$ 的幂就叫幂函数。}}
\begin{center}\begin{tikzpicture}[font=\small,>=stealth,line join=round]
\draw[gray!40,rounded corners=3pt] (0,0) rectangle (14.2,8.3);\node at (7.1,7.700000000000001) {固定指数2，改变系数会改变每一点的高度};
\begin{scope}[shift={(4.1,2.11)},x=1.6666666666666665cm,y=0.9199999999999999cm]\draw[->,gray] (-1.8,0)--(1.8,0) node[right]{$x$};\draw[->,gray] (0,-0.5)--(0,4.5) node[above]{$y$};\draw[blue,thick,solid,domain=-1.35:1.35,samples=180] plot(\x,{(\x)^2});\draw[teal,thick,dashed,domain=-1.35:1.35,samples=180] plot(\x,{2*(\x)^2});\draw[violet,thick,densely dotted,domain=-1.35:1.35,samples=180] plot(\x,{.5*(\x)^2});\fill[orange!85!black] (1,1) circle(2.2pt);\fill[orange!85!black] (1,2) circle(2.2pt);\fill[orange!85!black] (1,0.5) circle(2.2pt);\draw[gray,dashed] (1,0)--(1,3.9);\node[below] at (1,0){$1$};\end{scope}\node[anchor=west,align=left,text width=5.0cm] at (8.4,4.0) {蓝：$x^2$\\[9pt]青：$2x^2$\\[9pt]紫：$\frac12x^2$\\[9pt]$x=1$处高度为$1,2,\frac12$\\[9pt]严格定义只取系数1};
\node[text width=12.6cm,align=center] at (7.1,.65) {同一指数也能有不同高度；固定系数1，才是单独研究指数的基本幂函数。};
\end{tikzpicture}\end{center}


      $2x^2$、$-x^3$ 都是合法的函数，但不是高中严格定义下的幂函数。若题目明确写 $Ax^\alpha$ 是幂函数，先令 $A=1$，而不是 $A\ne0$。
    
\end{trapbox}
\clearpage
\begin{trapbox}
\textbf{\textcolor{CTrap}{易错二：认为系数为 $1$ 是一条需要证明的定理。}}
\begin{center}\begin{tikzpicture}[font=\small,>=stealth,line join=round]
\draw[gray!40,rounded corners=3pt] (0,0) rectangle (14.2,8.3);\node at (7.1,7.700000000000001) {易错二：只经过(1,1)，还不能判定是幂函数};
\begin{scope}[shift={(1.8741935483870968,2.9045454545454543)},x=1.9354838709677418cm,y=0.8363636363636363cm]\draw[->,gray] (-0.4,0)--(2.7,0) node[right]{$x$};\draw[->,gray] (0,-1.5)--(0,4) node[above]{$y$};\draw[teal,thick,solid,domain=0:2.3,samples=180] plot(\x,{\x});\draw[orange!85!black,thick,dashed,domain=0:2.3,samples=180] plot(\x,{2*\x-1});\fill[orange!85!black] (1,1) circle(2.2pt);\end{scope}\node[anchor=west,align=left,text width=5.0cm] at (8.4,4.0) {青：$y=x$是幂函数\\[9pt]橙：$y=2x-1$不是\\[9pt]两者都过$(1,1)$\\[9pt]但橙线有常数项$-1$\\[9pt]必须核对整个函数式};
\node[text width=12.6cm,align=center] at (7.1,.65) {图中的反例说明：过公共点是必要条件，不能单独当作判定任意函数为幂函数的充分条件。};
\end{tikzpicture}\end{center}


      这是定义规定。我们可以解释选择 $1$ 的好处，也可以从必过 $(1,1)$ 得到参数方程，但不能把人为定义当作自然推导出的定理。仅过 $(1,1)$ 也不能判定任意函数都是幂函数，例如 $y=2x-1$。
    
\end{trapbox}
\clearpage
\begin{trapbox}
\textbf{\textcolor{CTrap}{易错三：认为所有幂函数都经过原点。}}
\begin{center}\begin{tikzpicture}[font=\small,>=stealth,line join=round]
\draw[gray!40,rounded corners=3pt] (0,0) rectangle (14.2,8.8);\node at (7.1,8.200000000000001) {易错三：经过原点，与在(0,1)有空缺是两件事};
\node at (3.1500000000000004,7.3) {$y=x^2$};\begin{scope}[shift={(3.1500000000000004,3.0428571428571427)},x=1.225cm,y=1.0857142857142856cm]\draw[->,gray] (-2,0)--(2,0) node[right]{$x$};\draw[->,gray] (0,-0.5)--(0,3) node[above]{$y$};\draw[blue,thick,solid,domain=-1.5:1.5,samples=180] plot(\x,{(\x)^2});\fill[orange!85!black] (0,0) circle(2pt);\end{scope}\node[align=center,text width=5.8cm] at (3.1500000000000004,1.7) {正指数的完整图像\\包含$(0,0)$};\node at (9.95,7.3) {$y=x^0$};\begin{scope}[shift={(9.95,3.0428571428571427)},x=1.225cm,y=1.0857142857142856cm]\draw[->,gray] (-2,0)--(2,0) node[right]{$x$};\draw[->,gray] (0,-0.5)--(0,3) node[above]{$y$};\draw[gray,thick,solid,domain=-1.6:-0.02,samples=180] plot(\x,{1});\draw[gray,thick,solid,domain=0.02:1.6,samples=180] plot(\x,{1});\draw[gray,fill=white,thick] (0,1) circle(3pt);\end{scope}\node[align=center,text width=5.8cm] at (9.95,1.7) {$x\ne0$时为1\\空缺位置是$(0,1)$};
\node[text width=12.6cm,align=center] at (7.1,.65) {两图同尺度；右图的空缺在水平线$y=1$上的$(0,1)$。};
\end{tikzpicture}\end{center}


      只有 $\alpha>0$ 时，完整图像才经过 $(0,0)$。$\alpha=0$ 时，$0^0$ 在这里无意义，$y=x^0=1$ 在 $x=0$ 处缺少的是 $(0,1)$，而不是原点。$\alpha<0$ 时，零不能作为底数。第一象限的讨论本身不包含任何坐标轴上的点。
    
\end{trapbox}
\clearpage
\begin{trapbox}
\textbf{\textcolor{CTrap}{易错四：一律认为指数越大，幂值越大。}}
\begin{center}\begin{tikzpicture}[font=\small,>=stealth,line join=round]
\draw[gray!40,rounded corners=3pt] (0,0) rectangle (14.2,8.3);\node at (7.1,7.700000000000001) {同过(1,1)，底数跨过1时，上下顺序互换};
\begin{scope}[shift={(1.6625,2.068181818181818)},x=1.8750000000000002cm,y=0.8363636363636363cm]\draw[->,gray] (-0.3,0)--(2.9,0) node[right]{$x$};\draw[->,gray] (0,-0.5)--(0,5) node[above]{$y$};\draw[blue,thick,solid,domain=0.03:2.15,samples=180] plot(\x,{(\x)^2});\draw[teal,thick,dashed,domain=0.03:2.6,samples=180] plot(\x,{\x});\draw[violet,thick,densely dotted,domain=0.03:2.6,samples=180] plot(\x,{sqrt(\x)});\fill[orange!85!black] (1,1) circle(2.2pt);\draw[gray,dashed] (1,0)--(1,4.6);\node[below=5pt] at (1,0){$1$};\end{scope}\node[anchor=west,align=left,text width=5.0cm] at (8.4,4.0) {蓝实线：$x^2$\\[8pt]青虚线：$x$\\[8pt]紫点线：$\sqrt{x}$\\[8pt]$0<x<1$：$x^2<x<\sqrt{x}$\\[8pt]$x>1$：$x^2>x>\sqrt{x}$};
\node[text width=12.6cm,align=center] at (7.1,.65) {这里始终固定系数为1；三条曲线在正半轴上只有公共点(1,1)。};
\end{tikzpicture}\end{center}


      例如 $2^3>2^2$，却有 $(1/2)^3<(1/2)^2$。先判断底数大于 $1$、等于 $1$，还是在 $0$ 与 $1$ 之间；底数为 $1$ 时幂值全都相等。
    
\end{trapbox}
\clearpage
\begin{trapbox}
\textbf{\textcolor{CTrap}{易错五：混淆比较底数和比较指数。}}
\begin{center}\begin{tikzpicture}[font=\small,>=stealth,line join=round]
\draw[gray!40,rounded corners=3pt] (0,0) rectangle (14.2,8.8);\node at (7.1,8.200000000000001) {易错五：左图比较底数，右图比较指数};
\node at (3.1500000000000004,7.3) {$y=x^{-2}$};\begin{scope}[shift={(2.8,3.366666666666667)},x=0.7000000000000001cm,y=1.8333333333333335cm]\draw[->,gray] (-3,0)--(4,0) node[right]{$x$};\draw[->,gray] (0,-0.2)--(0,1.6) node[above]{$y$};\draw[blue,thick,solid,domain=0.85:3.5,samples=180] plot(\x,{1/((\x)^2)});\fill[orange!85!black] (2,0.25) circle(2pt);\fill[orange!85!black] (3,0.1111111111111111) circle(2pt);\node[below=4pt] at (2,0){$2$};\node[below=4pt] at (3,0){$3$};\end{scope}\node[align=center,text width=5.8cm] at (3.1500000000000004,1.7) {$2<3$，但$2^{-2}>3^{-2}$\\固定负指数，大小反向};\node at (9.95,7.3) {$y=2^\alpha$};\begin{scope}[shift={(9.6,3.366666666666667)},x=0.7000000000000001cm,y=1.8333333333333335cm]\draw[->,gray] (-3,0)--(4,0) node[right]{$\alpha$};\draw[->,gray] (0,-0.2)--(0,1.6) node[above]{$y$};\draw[orange!85!black,thick,solid,domain=-2.8:0.5,samples=180] plot(\x,{2^\x});\fill[orange!85!black] (-2,0.25) circle(2pt);\fill[orange!85!black] (-1,0.5) circle(2pt);\node[below=4pt] at (-2,0){$-2$};\node[below=4pt] at (-1,0){$-1$};\end{scope}\node[align=center,text width=5.8cm] at (9.95,1.7) {$-2<-1$，且$2^{-2}<2^{-1}$\\固定底数2，大小同向};
\node[text width=12.6cm,align=center] at (7.1,.65) {两图同尺度，但横轴意义不同；先辨认改变的是底数还是指数。};
\end{tikzpicture}\end{center}


      $2^{-2}>3^{-2}$ 比较的是底数；$2^{-2}<2^{-1}$ 比较的是指数。前者看指数正负，后者看底数与 $1$ 的关系。两个条件都变时不能直接套用其中一条。
    
\end{trapbox}
\clearpage
\begin{trapbox}
\textbf{\textcolor{CTrap}{易错六：把“越来越平”误认为函数值一定越来越小。}}
\begin{center}\begin{tikzpicture}[font=\small,>=stealth,line join=round]
\draw[gray!40,rounded corners=3pt] (0,0) rectangle (14.2,8.8);\node at (7.1,8.200000000000001) {易错六：都越来越平，却一个上升、一个下降};
\node at (3.1500000000000004,7.3) {$y=\sqrt{x}$};\begin{scope}[shift={(1.1,2.88)},x=1.0cm,y=0.76cm]\draw[->,gray] (-0.4,0)--(4.5,0) node[right]{$x$};\draw[->,gray] (0,-0.5)--(0,4.5) node[above]{$y$};\draw[violet,thick,solid,domain=0.05:4,samples=180] plot(\x,{sqrt(\x)});\draw[orange!85!black,thick,dashed,domain=0.2:0.8,samples=180] plot(\x,{0.7071067811865476+(0.7071067811865476)*(\x-(0.5))});\draw[orange!85!black,thick,dashed,domain=1.5:2.5,samples=180] plot(\x,{1.4142135623730951+(0.3535533905932738)*(\x-(2))});\fill[orange!85!black] (0.5,0.7071067811865476) circle(2pt);\fill[orange!85!black] (2,1.4142135623730951) circle(2pt);\end{scope}\node[align=center,text width=5.8cm] at (3.1500000000000004,1.7) {切线斜率均为正\\由$\frac1{\sqrt2}$减到$\frac1{2\sqrt2}$};\node at (9.95,7.3) {$y=1/x$};\begin{scope}[shift={(7.9,2.88)},x=1.0cm,y=0.76cm]\draw[->,gray] (-0.4,0)--(4.5,0) node[right]{$x$};\draw[->,gray] (0,-0.5)--(0,4.5) node[above]{$y$};\draw[orange!85!black,thick,solid,domain=0.25:4,samples=180] plot(\x,{1/\x});\draw[orange!85!black,thick,dashed,domain=0.3:0.7,samples=180] plot(\x,{2.0+(-4.0)*(\x-(0.5))});\draw[orange!85!black,thick,dashed,domain=1.4:2.6,samples=180] plot(\x,{0.5+(-0.25)*(\x-(2))});\fill[orange!85!black] (0.5,2) circle(2pt);\fill[orange!85!black] (2,0.5) circle(2pt);\end{scope}\node[align=center,text width=5.8cm] at (9.95,1.7) {切线斜率均为负\\由$-4$变为$-\frac14$};
\node[text width=12.6cm,align=center] at (7.1,.65) {两图同尺度，切点横坐标均取0.5与2；“更平”看斜率绝对值，“升降”看正负。};
\end{tikzpicture}\end{center}


      $\sqrt{x}$ 上升越来越慢，$1/x$ 下降越来越慢，两者都越来越平。判断方向看斜率的正负，判断陡峭程度看斜率绝对值的大小。
    
\end{trapbox}
\clearpage
\begin{trapbox}
\textbf{\textcolor{CTrap}{易错七：把正半轴上的单调性直接推广到整个定义域。}}
\begin{center}\begin{tikzpicture}[font=\small,>=stealth,line join=round]
\draw[gray!40,rounded corners=3pt] (0,0) rectangle (14.2,8.8);\node at (7.1,8.200000000000001) {易错七：正半轴的规律，不能直接覆盖负半轴};
\node at (3.1500000000000004,7.3) {$y=x^2$};\begin{scope}[shift={(3.1500000000000004,2.8454545454545457)},x=0.8166666666666668cm,y=0.6909090909090909cm]\draw[->,gray] (-3,0)--(3,0) node[right]{$x$};\draw[->,gray] (0,-0.5)--(0,5) node[above]{$y$};\draw[blue,thick,solid,domain=-2.1:2.1,samples=180] plot(\x,{(\x)^2});\fill[orange!85!black] (0,0) circle(2pt);\fill[orange!85!black] (1,1) circle(2pt);\fill[orange!85!black] (-1,1) circle(2pt);\end{scope}\node[align=center,text width=5.8cm] at (3.1500000000000004,1.7) {负半轴递减\\正半轴递增};\node at (9.95,7.3) {$y=x^{-2}$};\begin{scope}[shift={(9.95,2.8454545454545457)},x=0.8166666666666668cm,y=0.6909090909090909cm]\draw[->,gray] (-3,0)--(3,0) node[right]{$x$};\draw[->,gray] (0,-0.5)--(0,5) node[above]{$y$};\draw[blue,thick,solid,domain=-2.7:-0.46,samples=180] plot(\x,{1/((\x)^2)});\draw[blue,thick,solid,domain=0.46:2.7,samples=180] plot(\x,{1/((\x)^2)});\fill[orange!85!black] (1,1) circle(2pt);\fill[orange!85!black] (-1,1) circle(2pt);\end{scope}\node[align=center,text width=5.8cm] at (9.95,1.7) {负半轴递增\\正半轴递减，0不在定义域};
\node[text width=12.6cm,align=center] at (7.1,.65) {两图同尺度；每个单调区间必须单独写出，不能跨过0把两支当成一段。};
\end{tikzpicture}\end{center}


      如 $x^{-2}$ 在 $(-\infty,0)$ 上递增，在 $(0,+\infty)$ 上递减；$x^2$ 在负半轴上递减，在正半轴上递增。分数指数还要检查根号是否有意义。遇到偶函数可先改用绝对值，但必须先检查定义域。
    
\end{trapbox}
\clearpage
\begin{trapbox}
\textbf{\textcolor{CTrap}{易错八：只解不等式，忘记排除无意义的点。}}
\begin{center}\begin{tikzpicture}[font=\small,>=stealth,line join=round]
\draw[gray!40,rounded corners=3pt] (0,0) rectangle (14.2,8.8);\node at (7.1,8.200000000000001) {例4与易错八：先比较距离，再删去无意义输入};
\begin{scope}[shift={(5,5.9)},x=1.45cm]\draw[->,gray] (-3.1,0)--(3.8,0);\fill[blue] (-1,0) circle(2pt);\fill[blue] (3,0) circle(2pt);\fill[orange!85!black] (0,0) circle(2pt);\node[below=5pt] at (-1,0){$-1$};\node[below=5pt] at (3,0){$3$};\node[below=5pt] at (0,0){$a=0$};\draw[<->,blue] (-1,.65)--(0,.65) node[midway,above]{$1$};\draw[<->,violet] (0,1.25)--(3,1.25) node[midway,above]{$3$};\end{scope}\node[align=center] at (7.1,4.4){$|a+1|<|3-a|$：a离$-1$比离3更近，中点为1};\begin{scope}[shift={(5,2.8)},x=1.45cm]\draw[->,gray] (-3.1,0)--(3.8,0);\draw[blue,line width=2pt] (-3,0)--(-1,0);\draw[blue,line width=2pt] (-1,0)--(1,0);\foreach \v in {-1,1}{\draw[blue,fill=white,thick] (\v,0) circle(3pt);\node[below=6pt] at (\v,0){$\v$};}\draw[orange!85!black,fill=white] (3,0) circle(3pt);\node[below=6pt] at (3,0){$3$};\draw[<- ,blue,thick] (-3,0)--(-2.8,0);\end{scope}\node at (7.1,1.5){解集：$(-\infty,-1)\cup(-1,1)$};
\node[text width=12.6cm,align=center] at (7.1,.65) {a=-1与a=3使某个函数输入为0；求得a<1后，仍须去掉其中的a=-1。};
\end{tikzpicture}\end{center}


      对 $f(x)=1/|x|^{2/3}$，$f(a+1)>f(3-a)$ 虽可化为 $a<1$，但必须排除 $a=-1$ 和 $a=3$；最终只需在所得范围内去掉 $-1$。
    
\end{trapbox}
\clearpage
\begin{trapbox}
\textbf{\textcolor{CTrap}{易错九：把系数变化后的性质照搬给基本幂函数，或反过来。}}
\begin{center}\begin{tikzpicture}[font=\small,>=stealth,line join=round]
\draw[gray!40,rounded corners=3pt] (0,0) rectangle (14.2,8.3);\node at (7.1,7.700000000000001) {同样平方，负系数把曲线翻到横轴下方};
\begin{scope}[shift={(2.1,4.46111111111111)},x=2.5cm,y=0.5111111111111111cm]\draw[->,gray] (-0.4,0)--(2,0) node[right]{$x$};\draw[->,gray] (0,-5.5)--(0,3.5) node[above]{$y$};\draw[blue,thick,solid,domain=0.03:1.5,samples=180] plot(\x,{(\x)^2});\draw[orange!85!black,thick,dashed,domain=0.03:1.5,samples=180] plot(\x,{-2*(\x)^2});\fill[orange!85!black] (1,1) circle(2.2pt);\fill[orange!85!black] (1,-2) circle(2.2pt);\node[above left=4pt] at (1,1){$(1,1)$};\node[below left=4pt] at (1,-2){$(1,-2)$};\end{scope}\node[anchor=west,align=left,text width=5.0cm] at (8.4,4.0) {蓝：$y=x^2$\\[9pt]正半轴严格递增\\[9pt]橙：$y=-2x^2$\\[9pt]正半轴严格递减\\[9pt]高度先乘2，再变号};
\node[text width=12.6cm,align=center] at (7.1,.65) {曲线按真实函数绘制；系数改变后的函数不能直接照搬基本幂函数的五类结论。};
\end{tikzpicture}\end{center}


      $x^2$ 在正半轴上递增，$-2x^2$ 在正半轴上递减。正系数保持增减方向，负系数颠倒方向，零系数使函数在原式有意义的范围内恒为 $0$。本结论的五种走势专门针对系数为 $1$ 的情形。
  
\end{trapbox}
